\section{The degree-eight arithmetic line: primitive image and the missing second coordinate}

We retain $K=\mathbf Q(i)$, $E:y^2=x^3-x$, and the prime $5$.
The free graded Lie algebra on the chosen CM homology basis $x,y$
is embedded in its associative tensor algebra. Derivation degree $m$
means that $x,y$ are sent to Lie words of length $m+1$.
The tensor calculations are exact over $\mathbf Q$, with their
lattices localized at $5$. G.9 additionally computes an actual
arithmetic image by a finite Kummer certificate and integral
Johnson-filtration operations.

\begin{statement}{G.1}{Open}{The arithmetic scalar and Lead 15}
The scalar $\lambda$ in the actual arithmetic line
\[
 \frac{d^3\chi^{(4)}_4}{267840}(A_8+\lambda B_8)
\]
is not evaluated here. Lead 15, which asserts $\lambda\in\mathbf Q$,
also remains open. G.4 independently verifies the two ambient vectors
and their coordinates; G.5 verifies that every $5$-adic scalar is
compatible with the universal elliptic and finite CM constraints.
In particular those constraints admit
$\lambda=31/20+\sqrt6\in\mathbf Q_5$
and do not imply rationality. This is a counterexample to that proposed
method of deducing rationality, not a counterexample for the actual
arithmetic scalar.
G.9 proves the genuine arithmetic constraint
$\lambda-31/20\in\mathbf Z_5$; in particular
$v_5(\lambda)=-1$ and $5\lambda\equiv4\pmod5$.
The same irrational example also satisfies this congruence.

The motivic argument in Lead 15 would need a rational subspace in the
\emph{specified} degree-eight derivation space, together with a
comparison identifying its scalar extension with the actual Galois
line. Merely having Betti and de Rham rational structures does not
give that comparison or identify the two structures in these
coordinates. No such theorem is supplied or proved in this return.
\end{statement}

\begin{statement}{G.2}{Open}{The actual integral arithmetic image}
The values of the second coordinate $\nu$ on an actual element of
the eighth Johnson subgroup of $\Delta$ are not calculated here.
G.9 determines the first-coordinate ideal to be $\mathbf Z_5$ and
proves that the arithmetic image has index $1$ in the saturation
of its rational line. More precisely it is
$\mathbf Z_5(C_8+\delta B_8)$ for the still-unknown
$\delta=\lambda-31/20\in\mathbf Z_5$. There exists an actual
eighth-Johnson element with coordinates $(\alpha,\nu)=(1,\delta)$,
but the second value has not been measured. Thus the requested
explicit embedded generator remains open. The rank-one image
has infinite index in the rank-two ambient plane.
\end{statement}

\begin{statement}{G.3}{Open}{The repeated-character arithmetic embeddings through degree fourteen}
No complete arithmetic generator list or integral image lattice in
all repeated-character spaces of degrees eight through fourteen is
obtained. The degree-eight measurement still missing in G.1 is
already part of that requested list. A family of universal
derivations, a metabelian character, or dimensions of the lower
bracket images does not determine those generator vectors. The
calculations below are therefore recorded as ambient certificates,
with no promotion to a description of the Galois subgroup.
\end{statement}

\begin{statement}{G.4}{Proved}{An independent exact certificate for the degree-eight plane}
Let $F(x)=y$, $F(y)=0$, and for even $m\geq2$ put
\[
 \epsilon_m(x)=\operatorname{ad}_x^m(y),\qquad
 \epsilon_m(y)=\sum_{j=0}^{m/2-1}(-1)^j
 [\operatorname{ad}_x^j(y),\operatorname{ad}_x^{m-1-j}(y)].
\]
Set
\[
 A_8=\frac{\operatorname{ad}_F^3\epsilon_8}{120},\quad
 B_8=\frac{[\epsilon_4,\operatorname{ad}_F^2\epsilon_4]}2,
 \quad C_8=A_8+\frac{31}{20}B_8.
\]
These vectors annihilate $[x,y]$, have weight $(4,4)$, and change
sign under $x\leftrightarrow y$. For the integral word-coordinate
functionals
\[
 \alpha(D)=-[x^5y^4]D(x),\qquad
 \nu(D)=[x^4yxy^3]D(x)+4[x^5y^4]D(x),
\]
one has
\[
 \begin{array}{c|rrr}
 &A_8&B_8&C_8\\\hline
 \alpha&1&0&1\\
 \nu&-31/20&1&0
 \end{array}.
\]
Writing $W_8=\mathbf Q_5A_8\oplus\mathbf Q_5B_8$, its saturated
lattice of integral Lie derivations is
\[
 W_8\cap\operatorname{Hom}_{\mathbf Z_5}
       (\mathbf Z_5x\oplus\mathbf Z_5y,L_9(\mathbf Z_5x\oplus\mathbf Z_5y))
       =\mathbf Z_5C_8\oplus\mathbf Z_5B_8.
\]
In the free metabelian Lie algebra, written
$\mathbf Qx\oplus\mathbf Qy\oplus\mathbf Q[u,v][x,y]$ with
$u=\operatorname{ad}_x$, $v=\operatorname{ad}_y$ on the commutator
module, the values of $A_8$ are
$u^4v^3[x,y]$ and $u^3v^4[x,y]$, while both values of $B_8$ vanish.
\end{statement}
\begin{proof}
The executable certificate \texttt{code/fgh\_degree8.py} starts from
these formulas, with its own tensor multiplication, commutator, and
derivation implementation; it imports no attached calculation.
It computes all rational associative coefficients, checks the
special equation and the two symmetries, and decomposes each value
in the standard Lyndon basis by triangular subtraction. At each
subtraction the leading word is checked to be Lyndon and its
standard bracket has leading coefficient $1$. The complete vectors
are recorded in \texttt{code/fgh\_degree8.json}.

For the two selected words its coefficients are
\[
 \begin{array}{c|rrr}
 &A_8(x)&B_8(x)&C_8(x)\\\hline
 {[x^5y^4]}&-1&0&-1\\
 {[x^4yxy^3]}&49/20&1&4
 \end{array},
\]
which gives the displayed matrix. The Lyndon coefficients of $B_8$
are integers; those of $C_8$ have denominators only $1,2,4$.
Thus both are integral at $5$. Conversely, an integral vector
$D\in W_8$ has integral word coefficients, so
$\alpha(D),\nu(D)\in\mathbf Z_5$. The coordinate matrix gives
$D=\alpha(D)C_8+\nu(D)B_8$. This proves saturation, without
identifying any lattice vector as arithmetic.

For an independent check of the metabelian assertion, the program
evaluates each Lyndon bracket in the explicit metabelian Lie
algebra above. The bracket of
$ax+by+p(u,v)[x,y]$ and $cx+dy+q(u,v)[x,y]$ has commutator
coordinate
\[
 ad-bc+(au+bv)q-(cu+dv)p.
\]
This elementary formula supplies the stated result and directly
checks that the second direction is invisible metabelianly.
\end{proof}

\begin{statement}{G.5}{Proved}{Universal equations do not select a rational direction}
For every $\lambda\in\mathbf Q_5$, the line
$\mathbf Q_5(A_8+\lambda B_8)$ exponentiates in the completed
positive geometric subgroup of Enriquez's graded universal elliptic
kernel. It has CM character $(4,4)$ and the prescribed interchange
sign. This remains true for an irrational scalar: the root of
$T^2-6$ congruent to $1$ modulo $5$ belongs to $\mathbf Q_5$ and
is not rational. A line in this plane is defined over $\mathbf Q$
if and only if its normalized scalar $\lambda$ is rational.
\end{statement}
\begin{proof}
Enriquez, Proposition 4.19, equation (48), defines
$\delta_{2n}=\epsilon_{2n+2}$ and places it in
$\mathfrak r_{\mathrm{ell}}^{\mathrm{gr}}$.
The paragraph following equation (49) gives
\[
 \exp(\widehat{\mathfrak b}_3^+)\rtimes\mathrm{SL}_2
       \subset R_{\mathrm{ell}}^{\mathrm{gr}},\qquad
 \mathfrak b_3=\langle\mathfrak{sl}_2,\delta_{2n}\rangle
       .
\]
The inclusion is the one explicitly stated there \cite{fghEnriquez}.
The half-sum in (48) pairs into the formula for $\epsilon_m(y)$
in G.4. The source uses the opposite derivation bracket, which does
not change the generated Lie subalgebra. Both $A_8$ and $B_8$
therefore belong to its positive part. Exponentiation of any scalar
linear combination converges in the degree filtration and gives
the asserted group elements to all orders. G.4 checks the CM weight
and interchange sign exactly. The scalar $\sqrt6$ exists by Hensel's
lemma since $1^2-6\equiv0\pmod5$ and $2$ is a unit, and is
irrational since $6$ is not a rational square.

Finally $A_8,B_8$ are rational and independent. A rational vector
spanning the normalized line has a nonzero $A_8$ coefficient and
therefore a rational ratio of its $B_8$ and $A_8$ coefficients.
That ratio is $\lambda$. The converse is immediate. No arithmetic
realization of the irrational candidate is asserted.
\end{proof}

\begin{statement}{G.6}{Proved}{The exact conditional lattice and index formula}
Let $L=\mathbf Q_5(A_8+\lambda B_8)$ be any line in this plane,
put $\delta=\lambda-31/20$, and set
\[
 m=\max\{0,-v_5(\delta)\},\qquad v_5(0)=+\infty.
\]
Then
\[
 L\cap(\mathbf Z_5C_8\oplus\mathbf Z_5B_8)
     =5^m\mathbf Z_5(C_8+\delta B_8).
\]
If $\mathcal A\ne0$ is the compact image of a continuous
homomorphism from a profinite group to this integral line and
$\alpha(\mathcal A)=5^a\mathbf Z_5$, then $a\geq m$ and
\[
 \mathcal A=5^a\mathbf Z_5(C_8+\delta B_8),\qquad
 [L\cap\Lambda:\mathcal A]=5^{a-m},
 \quad\Lambda=\mathbf Z_5C_8\oplus\mathbf Z_5B_8.
\]
Thus this formula applies to the character-$(4,4)$ image of an
eighth Johnson homomorphism whenever its rational line is the
one specified in the prompt. G.9 proves $a=m=0$ for that actual
homomorphism; the value of $\delta$ still has to be measured.
\end{statement}
\begin{proof}
The coordinates of $t(C_8+\delta B_8)$ are $(t,t\delta)$.
Both are integral exactly when $v_5(t)\geq0$ and
$v_5(t)+v_5(\delta)\geq0$, giving the first formula.
A compact additive subgroup of $\mathbf Q_5$ is closed and stable
under $\mathbf Z_5$, by taking limits of multiplication by integers.
Its nonzero integral image is consequently an ideal
$5^a\mathbf Z_5$. Since $\alpha$ is injective on $L$, it determines
the whole subgroup and the index. The full rank-two quotient
$\Lambda/\mathcal A$ also has a rank-one part, so its cardinal
index is infinite; the finite displayed index is taken only inside
the saturation of $L$.
\end{proof}

\begin{statement}{G.7}{Proved}{A finite nonmetabelian measurement that would decide the scalar}
Let $P$ be the free pro-$5$ group on lifts $a,b$ of the integral
homology basis. An actual automorphism $\sigma$ in its eighth
Johnson subgroup determines $D_8(\sigma)$ by the leading terms of
$\sigma(a)a^{-1}$ and $\sigma(b)b^{-1}$. For every $r\geq1$,
the two coordinates $\alpha,\nu$ modulo $5^r$ can be read from its
action on the characteristic finite quotient supplied by
\[
 (\mathbf Z/5^r)\langle X,Y\rangle/(X,Y)^{10},
 \qquad a\mapsto1+X,\quad b\mapsto1+Y.
\]
On a vector of the arithmetic line with $\alpha\ne0$,
\[
       \lambda=\frac{31}{20}+\frac{\nu}{\alpha}.
\]
This is an explicit test on a genuine arithmetic action; it does
not provide such an action or its measured values.
\end{statement}
\begin{proof}
The completed group algebra of a free pro-$5$ group is the completed
noncommutative power-series algebra under these substitutions.
Its augmentation ideal and the ideal generated by $5^r$ are
characteristic. The quotient algebra is finite and its augmentation
units form a finite $5$-group, as follows from the finite additive
successive powers of the augmentation ideal. Hence the indicated
group quotient and automorphism action are defined.

For an eighth Johnson element, all positive terms of lengths less
than nine in $\sigma(a)a^{-1}-1$ vanish. The degree-nine term is
the Lie element $D_8(\sigma)(x)$ because leading group commutators
expand to the corresponding Lie commutators. The two integral
linear combinations of word coefficients in G.4 therefore compute
$\alpha,\nu$ modulo $5^r$ without division or a logarithm.
The ratio formula is the coordinate identity of G.6. In particular
metabelian data, which kill $B_8$, cannot replace this second
measurement.
\end{proof}

\begin{statement}{G.8}{Known}{The realization hypothesis in the proposed motivic shortcut}
Malkin's \emph{Motivic fundamental groups of CM elliptic curves
and geometry of Bianchi hyperbolic threefolds}, Theorem 11,
proves a depth-two symmetry for Hodge correlators; Theorem 16
constructs the corresponding Bianchi-complex map in that realization.
The transfer discussed immediately before Lemma 3, on printed
pages 12--13, invokes hypothetical regulator injectivity and the
motivic formalism. It therefore supplies no unconditional theorem
identifying the degree-eight $5$-adic arithmetic line with a rational
line in the fixed derivation coordinates used here
\cite{fghMalkin}.
\end{statement}
The primary text has been opened at those statements. This is a
precise limit on what that proposed source proves, not a proof that
the desired comparison cannot be established by another argument.

We also checked a more direct mixed-Tate route to Lead 15.
The subsection \emph{Motivic Sheaves} of
\cite{fghUniversalMEM} explicitly leaves the motivic
specialization of its compatible realizations as a question.
The fully faithful comparison in
\cite[Proposition 10.3.3]{RootHuberMuellerStach2017} has
subquotient-closed image; it does not identify an arbitrary
Tate-filtered Nori extension with a standard mixed Tate
motive. The equivalence in
\cite[Theorem 1.5]{RootTubach2025} assumes a motivic
\(t\)-structure. Finally,
\cite[Theorem A]{RootKupersRudenkoSierra2026} concerns
the standard mixed-Tate category over a number field,
and supplies no elliptic specialization comparison.
Thus the missing step in this route is an actual
construction of the balanced elliptic extension in that
category, with its two coordinate maps. Its one-dimensional
Tate-weight-four extension space would then control a
rational slope, but the known one-dimensional
\(5\)-adic image does not furnish that construction.
The precise attempted reduction and source checks are
recorded in \path{notes/root_motivic_rationality.md}.

\begin{statement}{G.9}{Proved}{The actual first-coordinate ideal, saturation index, and scalar congruence}
Let $J_m$ be the subgroup of $\Delta$ whose actual tangential
action on $P/\gamma_{m+1}P$ is trivial. In the precise
conductor-four normalization of Ishii--Goto, one has
\[
 \chi^{(4)}_4(\Delta)=\chi^{(4)}_4(J_8)=5\mathbf Z_5,
 \qquad
 \alpha\bigl(j_8(J_8)^{(4,4)}\bigr)=\mathbf Z_5.
\]
Here the superscript denotes the projection to CM character
$(4,4)$, and $\alpha$ is the integral coordinate of G.4.
Consequently the actual arithmetic line of G.1 has image
\[
 \mathcal A=\mathbf Z_5(C_8+\delta B_8),\qquad
 \delta=\lambda-\frac{31}{20}\in\mathbf Z_5,
 \qquad [L\cap\Lambda:\mathcal A]=1,
\]
where $L=\mathbf Q_5(A_8+\lambda B_8)$ and
$\Lambda=\mathbf Z_5C_8+\mathbf Z_5B_8$. In particular
\[
 v_5(\lambda)=-1,\qquad 5\lambda\equiv4\pmod5.
\]
An actual element of $J_8$ has $(4,4)$ coordinates $(1,\delta)$.
This computes its first coordinate and the saturation index,
but leaves its second coordinate $\delta$ unevaluated.
\end{statement}
\begin{proof}
\emph{The exact Kummer normalization and a primitive value.}
Write $\chi=\chi_5$ and $k_4=\chi^{(4)}_4$.
Definition 2.14 and Remark 2.15(2) of
\citet{fghIshiiGoto} define an integral cocycle class
in $H^1(G_K,\mathbf Z_5(4))$. On
$G_{K(\mu_{5^\infty})}$, its reduction modulo $5$ is the
Kummer character of
\[
 u_1=\prod_{a=1}^4(1-i\zeta_5^{\,4a})^{a^3}
       \quad\text{in }F^\times/F^{\times5},
 \qquad F=\mathbf Q(i,\zeta_5).
\]
Indeed the norm in that definition is the identity for $K=\mathbf
Q(i)$, $\zeta_{20}=i\zeta_5^4$, and the automorphism indexed by
$a$ fixes $i$ and sends $\zeta_5$ to $\zeta_5^a$. Thus no
Euler factor or divided Soul\'e normalization is being inserted.

At the prime of $F$ above $101$ defined by $i\mapsto10$ and
$\zeta_5\mapsto95$, the four factors before exponentiation are
$70,40,45,61$. After exponentiation by $1,8,27,64$ they are
$70,88,96,80$, and hence
\[
 u_1\equiv97\pmod{101},\qquad
 97^{20}\equiv36\ne1\pmod{101},\qquad95^{20}\equiv1\pmod{101}.
\]
The elements involved are units at this prime. The map
$v\mapsto\overline v^{\,20}$ kills fifth powers and the line
generated by $\zeta_5$, but not $u_1$. Therefore the Kummer
class of $u_1$ lies outside that cyclotomic line.
The extension $K(\mu_{5^\infty})/F$ has group $\mathbf Z_5$
and unique degree-five subextension $F(\zeta_{25})$,
whose Kummer line is generated by $\zeta_5$.
It follows that $k_4$ is nonzero modulo $5$ on
$G_{K(\mu_{5^\infty})}$. Its compact image is consequently
$\mathbf Z_5$.

\emph{Restriction to the full commutator subgroup.}
Modulo $5$, the coefficient action $\chi^4$ is trivial. Any
integral cocycle representative of $k_4$ therefore reduces
to an additive homomorphism on $G_K$ and vanishes on
its closed commutator subgroup $\Delta$. This proves
$k_4(\Delta)\subset5\mathbf Z_5$.
Choose $h\in G_{K(\mu_{5^\infty})}$ with $k_4(h)$ a unit.
Since $K\cap\mathbf Q(\mu_{5^\infty})=\mathbf Q$, choose
$g\in G_K$ with $\chi(g)=6$. The cocycle identity gives
\[
 k_4([g,h])=(6^4-1)k_4(h)=1295\,k_4(h).
\]
This value has valuation exactly one and $[g,h]\in\Delta$.
The restriction of $k_4$ to $\Delta$ is a continuous
homomorphism with compact image, so that image is
exactly $5\mathbf Z_5$.

\emph{An integral filter through the lower Johnson degrees.}
We give group operations, without dividing by a logarithm
or passing to a rational Zariski closure.
For $r\in\mathbf Z_5$, let $\widehat r\in\widehat{\mathbf Z}$
have $5$-component $r$ and every other prime component zero.
For $u\in\Delta$ its profinite power $u^{\widehat r}$
exists in the closed procyclic subgroup it generates.
Both its Johnson coordinates and its $k_4$ value are
multiplied by $r$.

Choose $k\in G_K$ with $\chi(k)\equiv2\pmod5$, and put
$T=\operatorname{Ad}(\rho_s(k))$ on each Johnson module.
For $u\in J_m$ and $r\in\mathbf Z_5$ define
\[
 \mathcal F_r(u)=kuk^{-1}u^{-\widehat r}.
\]
Normality and the Johnson homomorphism give
\[
 j_m(\mathcal F_r(u))=(T-r)j_m(u),\qquad
 k_4(\mathcal F_r(u))=(\chi(k)^4-r)\,k_4(u).
\]
The second identity uses $\chi|_\Delta=1$ and holds
exactly, even though $\mathcal F_r$ need not be a group
homomorphism.

The CM lift $J$ of $[i]$ commutes with the actual action:
on the tangential tower $t_n^{5^n}=t$ it is $t_n\mapsto i t_n$,
which is defined over $K$ and compatible with the tower.
Thus the allowed weight $(a,b)$ has $a-b\equiv0\pmod4$.
In particular odd Johnson images vanish integrally.
The complete list of allowed special characters below
degree eight is
\[
 \begin{array}{c|c|c}
  m& (a,b)\text{ and multiplicities}
        &\psi(k)^a\overline\psi(k)^b\bmod5\\ \hline
  2&(1,1):1&2\\
  4&(2,2):1&4\\
  6&(5,1):1,\ (3,3):2,\ (1,5):1&2,\ 3,\ 2
 \end{array}.
\]
For completeness, this list follows from the exact special
kernel
$L_{m+1}^2\to L_{m+2}$, $(u,v)\mapsto[u,y]+[x,v]$.
Its character multiplicity is
$\ell(a+1,b)+\ell(a,b+1)-\ell(a+1,b+1)$, where
$\ell$ is the two-variable Witt rank. Surjectivity follows
by Jacobi from generation by $x,y$. All character spaces
are integral direct summands in the CM eigenbasis.

We have $J_1=\Delta$ and $J_1=J_2$ by odd-degree vanishing.
Apply $\mathcal F_{\chi(k)}$ to kill degree two, then
$\mathcal F_{\chi(k)^2}$ to kill degree four.
At degree six, apply in succession the three operators
with
\[
 r=\psi(k)^5\overline\psi(k),\quad
   \chi(k)^3,\quad
   \psi(k)\overline\psi(k)^5.
\]
Their product annihilates that whole Johnson module, including
the multiplicity-two middle space, since $T$ acts on each
character space by its displayed scalar. The intervening
odd degrees vanish, so the resulting element belongs to
$J_8$. Every multiplier $\chi(k)^4-r$ is a $5$-adic unit:
the target residue is $1$ and the five excluded residues
are $2,4,2,3,2$. Their product is $4$ modulo $5$.
Thus this finite sequence of actual group operations takes
the full $k_4$ image $5\mathbf Z_5$ into $k_4(J_8)$
without changing that ideal. The reverse inclusion was
already proved, giving $k_4(J_8)=5\mathbf Z_5$.

\emph{Independent comparison with the actual Magnus coordinate.}
Here are the source normalizations, including the factor
which cannot be read from the uncorrected 1995 formula.
Nakamura's (4.8) defines his metabelian power series
$\beta$ by
$\sigma(a_i)a_i^{-1}=(T_i\beta)\cdot z$ with
$z=[a_1,a_2]^{-1}$. Its leading Lie term is therefore
$-X_i\beta\,[x,y]$. In the notation
$D(x)=Xh[x,y]$, $D(y)=Yh[x,y]$, this means
$h=-\beta$. His Theorem A, equivalently Corollary 4.12,
gives
\[
 [X^3Y^3]h
       =-\frac{\kappa_{33}}{(1-5^6)(3!)^2}
       \qquad\text{on }J_8.
\]
These statements are checked directly in
\citet[Sections 4.3--4.9 and Corollary 4.12]{Nakamura1995}.
For an element of $J_8$, the genuine tangential lift
equals the canonical lift used there: their difference
is an inner power of the boundary, and its degree-two
term would be nonzero unless that power were zero.

The published correction to (3.11.4) in
\citet[p.~492, Errata to N95]{fghNakamura2013}
makes the raw theta-product Kummer exponent
$12\kappa_{33}$. The projector eigenbasis has columns
$(1,-i_5)/2,(1,i_5)/2$ and determinant $d=i_5/2$.
On replacing its second vector by $d^{-1}$ times that
vector, reindexing the exponents $a^3b^3$ multiplies
the raw coefficient by $d^3$. Accordingly
\[
 12\kappa_{33}=d^3\chi^{\rm ell}_{(4,4)}.
\]
Theorem 1.7 of \citet{fghIshiiGoto}, with
$w_K=4$, $B_4=-1/30$ and $p=5$, gives
\[
 \chi^{\rm ell}_{(4,4)}
 =-\frac{6\cdot4\,B_4}{4}
       \frac{1-5^6}{1-5^3}\,k_4
 =\frac{126}{5}\,k_4.
\]
Thus, using G.4's independently executed metabelian
extraction $h(A_8)=X^3Y^3$, $h(B_8)=0$, and its
word-coordinate identity $\alpha(A_8)=1$,
\[
 \alpha(j_8(\sigma)^{(4,4)})
 =-\frac{d^3(126/5)k_4(\sigma)}
          {12(1-5^6)(3!)^2}
 =\frac{d^3}{267840}\,k_4(\sigma).
\]
This checks the actual coefficient, rather than only
the normalization of an ambient derivation.
Since $d$ is a unit and $267840=5\cdot53568$ with
$53568\in\mathbf Z_5^\times$, the proved $k_4$ image
gives $\alpha(j_8(J_8)^{(4,4)})=\mathbf Z_5$.

\emph{The rational-line input for the lattice conclusion.}
The line used here is the proved rational-image assertion
LM8 G.15 and (G.25), not the still-open evaluation in G.1.
Its inputs can be checked separately.
Ishii's Theorem 1.2 and Lemma 1.5 give one outer
generator in each positive pair $a\equiv b\pmod4$
other than $(1,1)$, for $(K,p)=(\mathbf Q(i),5)$
\cite{fghIshiiGraded}. The ray field $K(5)=\mathbf
Q(\zeta_{20})$ has class number one: its Minkowski
bound is below $13$, its only primes of norm below
that bound are the two above $5$, and
$\operatorname{Norm}(\zeta_{20}^3-\zeta_{20}+1)=5$.
In $K(25)$ the ray group is
$(\mathbf Z/25)^{\times\,2}/\mu_4$ of order $100$.
At $(2+i)$, inertia has order $20$ and the image
of Frobenius in the order-five quotient is represented
by $9$, using $i\equiv7$ at the conjugate prime.
It generates that quotient; hence there are exactly
two primes above $5$, as required by Lemma 1.5.
Below degree eight the only outer degree-four
generator has character $(2,2)$, so its self-bracket
is zero. Thus the $(4,4)$ image has dimension at
most one. The nonzero coefficient just proved makes
its dimension exactly one. Genuine inner special
derivations contribute only in degree two:
$[w,[x,y]]=0$ forces $w$ to be a multiple of $[x,y]$.

For its placement in $W_8$, the universal monodromy
comparison of Hain--Matsumoto, Proposition 20.1,
Corollary 20.2 and Section 22, has constant generators
$\mathbf Q(2n+1)$ only, and geometric generators
the $\epsilon_{2n}$ modules
\cite{fghUniversalMEM}. In ordinary degree eight
there is no constant Tate-weight-four generator;
the degree-two boundary generator is central.
The geometric contributions are the $\epsilon_8$
module and the bracket of two degree-four modules.
Their $(4,4)$ parts are exactly the independently
checked vectors $A_8,B_8$ of G.4.
This gives the stated rank-one subspace of $W_8$
without asserting its unknown ratio is rational.
It does not transfer a formula at the Tate fibre
to the CM fibre.

Finally the projection to $(4,4)$ is integral: it
selects the appropriate multidegrees in the integral
free Lie algebra on the integral CM basis. Its image
therefore lies in $\Lambda$. In the actual rational
line of G.1 it has coordinates $(t,t\delta)$.
Taking $t=1$ proves $\delta\in\mathbf Z_5$.
G.6 now gives the exact image and saturation index.
The value $31/20$ has valuation $-1$, so adding the
integral $\delta$ gives the claimed valuation and
congruence of $\lambda$. The index in the full
rank-two lattice $\Lambda$ remains infinite.

The independent exact program
\texttt{code/fgh\_primitive\_degree8.py} checks the
residue certificate, all the lower character
multiplicities, every possible residual CM pair with
product $2$, and the numerical normalization.
Its JSON and stdout are retained. The first residue
certificate was also executed independently in
\texttt{code/de\_soule\_conductor4.py}.
\end{proof}
